\subsection{A finite bound for the manuscript character in every height}

\begin{lemma}[Comparison after a common good-reduction extension]
\label{lem:B-RD17-common}
Let $M,N$ be height-$h$ BT groups over $K$, each equipped with a
filtration by BT subgroups whose graded pieces have good reduction
over the same finite Galois extension $L/K$. Let $c$ be an integer
such that $p^c$ times every generic-fibre homomorphism between finite
flat commutative $p$-group schemes over $\mathcal O_L$ extends over
$\mathcal O_L$. If
\[
 M[p^n]\simeq N[p^n],\qquad n>2ch^2,
                                                        \tag{B-RD17.1}
\]
then their manuscript characters agree on all prime-to-$p$ inertia
elements. No good-reduction assertion about the extensions between
the graded pieces is assumed.
\end{lemma}
\begin{proof}
Use the given isomorphism to identify the two generic finite groups
with $X$. The two BT filtrations induce filtrations
\[
 0=X_0\subset\cdots\subset X_r=X,
 \qquad 0=Y_0\subset\cdots\subset Y_s=X.
\]
Write $h_i$ and $h'_j$ for the heights of their respective graded BT
groups. Thus $\sum_i h_i=\sum_j h'_j=h$.
The Zassenhaus refinement has common generic pieces
\[
 Q_{ij}=
 \frac{X_i\cap Y_j}
 {(X_{i-1}\cap Y_j)+(X_i\cap Y_{j-1})}.
                                                        \tag{B-RD17.2}
\]
In the good model of the $i$th graded group at level $p^n$, take
schematic closures of the subgroup filtration induced by the
$Y_j$. Consecutive quotients give finite flat models $G_{ij}$ of
$Q_{ij}$. Refine the good model of the $j$th graded group from $N$
in the same way, obtaining another model $H_{ij}$. These quotients
need not be truncated BT groups; only their finite flatness is used.
All closures and quotients are stable under the descent action.
The descent actions on the good BT models exist uniquely by Tate's
extension theorem for homomorphisms over a normal base.

Let $k_L$ be the residue field, and extend it algebraically if
necessary. The special Dieudonne module of $G_{ij}$ is a subquotient
of the free $W(k_L)/p^n$-module belonging to the $i$th graded BT
model. Consequently it can be generated by at most $h_i$ elements
over the DVR $W(k_L)$. Similarly the module for $H_{ij}$ has at most
$h'_j$ generators. Here exactness of Dieudonne theory over a perfect
field is applied only to special fibres of exact finite flat
sequences.

The generic identification of $G_{ij}$ and $H_{ij}$, and its inverse,
extend after multiplication by $p^c$. The two composites are
$p^{2c}$, because this equality holds generically and the source
coordinate algebras are flat. Uniqueness of these extensions makes
them compatible with inertia. Fix a prime-to-$p$ element $\gamma$
of inertia, of order $m$. After extending the perfect residue field
so that all $m$th roots of unity are present, the integral
idempotents
\[
 \frac1m\sum_{a=0}^{m-1}\zeta^{-a}\gamma^a
\]
split every special Dieudonne module into $\zeta$-eigensummands.
On a fixed eigensummand, the kernel of either comparison map is
annihilated by $p^{2c}$, and its cokernel is also annihilated by
$p^{2c}$. A module with at most $d$ generators over a DVR has
$p^{2c}$-torsion of length at most $2cd$, and any quotient killed
by $p^{2c}$ has length at most $2cd$. It follows that the two
$\zeta$-eigenlengths differ by at most
\[
 2c\max(h_i,h'_j)\leq2ch_i h'_j.
                                                        \tag{B-RD17.3}
\]
(The last inequality uses that all nonzero graded heights are
positive integers.)

Exactness along each refinement adds these lengths. A good BT group
of height $d$ has a free integral Dieudonne module of rank $d$, and
its $p^n$-torsion contributes $n$ times the multiplicity of each
$\gamma$-eigenvalue in that module. Thus, if $a_\zeta,b_\zeta$ are
the multiplicities in the sums of graded isocrystals for $M,N$,
\[
 \bigl|n(a_\zeta-b_\zeta)\bigr|
 \leq 2c\sum_{i,j}h_i h'_j=2ch^2.
\]
Under (B-RD17.1), the integer $a_\zeta-b_\zeta$ is zero. This holds
for every eigenvalue and every prime-to-$p$ inertia element, proving
the assertion. In particular, applying the same argument to two
filtrations of a single $M$, with arbitrarily large $n$, proves
independence of the prime-to-$p$ character from the chosen filtration.
\end{proof}

\begin{lemma}[A uniform ramification degree for the graded pieces]
\label{lem:B-RD17-degree}
Put $d=4h^2$ and
\[
 D_{p,h}=
 \begin{cases}
  |\GL_d(\mathbf F_p)|,&p>2,\\
  |\GL_d(\mathbf Z/4\mathbf Z)|
     =2^{d^2}|\GL_d(\mathbf F_2)|,&p=2.
 \end{cases}
\]
For two height-$h$ groups satisfying the monodromy condition, their
graded pieces have simultaneous good reduction over a finite Galois
extension $L/K$ with ramification index at most $D_{p,h}$.
\end{lemma}
\begin{proof}
Choose any finite Galois good-reduction extension $L_0/K$ for all
the graded pieces. Let $I_0$ be its inertia group. Their direct sum
has height $2h$; denote its special Dieudonne isocrystal by $D$.
The image $H$ of $I_0$ lies in
$\End_F(D)^\times$, where $F$ denotes semilinear Frobenius.
This endomorphism algebra has $\mathbf Q_p$-dimension at most
$(2h)^2=d$. Here is an elementary justification not requiring a
classification of isocrystals. In any semilinear Frobenius module,
linearly independent fixed vectors over the fixed field
$W(k_{L_0})[1/p]^{F}=\mathbf Q_p$ remain linearly independent over
$W(k_{L_0})[1/p]$: normalize a shortest relation to have one
coefficient $1$, apply Frobenius, and subtract. A shorter relation
would result unless every coefficient were fixed. Apply this to
conjugation by Frobenius on the endomorphism module, of rank $d$.

Left multiplication embeds $H$ in a general linear group over
$\mathbf Q_p$ of degree at most $d$. A finite group preserves a
$\mathbf Z_p$-lattice, obtained by summing the translates of any
lattice. For $p>2$, the kernel of reduction modulo $p$ in
$\GL_d(\mathbf Z_p)$ is torsion free. For $p=2$, the kernel of
reduction modulo $4$ is torsion free. The $p$-adic logarithm, or the
binomial valuation argument, proves both assertions. Thus
$|H|\leq D_{p,h}$.

Let $J$ be the kernel of $I_0\to H$. It is normal in
$\operatorname{Gal}(L_0/K)$, and $L=L_0^J$ is Galois over $K$ with
$e(L/K)=|H|$. Each rational graded Tate module is crystalline over
$L$. Indeed crystalline comparison identifies its special
isocrystal, with its descent action, with the crystalline period
module, up to the dual or Tate twist dictated by the convention.
The twist has trivial inertia action. The group $J$ acts trivially,
and $L_0,L$ have the same residue field. Hence
\[
 D_{\cris,L}(V)=D_{\cris,L_0}(V)^J
\]
has full dimension. The period-invariant and unramified-completion
statements used here are also explicitly treated in
\cite[\S\S5.1.5--5.1.7, pp.156--158]{Fontaine1994Semi}.
The required comparison, including functoriality,
is given in \cite[Lemma~2.2.4 and the proof of Corollary~2.2.6]{Kisin2006Crystalline}.
Corollary~2.2.6 of that paper, valid for arbitrary perfect residue
field and every prime, now gives a good BT group with this rational
Tate module.

For completeness, the actual graded lattice, rather than just its
rational representation, also has a good model. If $T_0$ is the Tate
lattice of one good model, rescale the desired lattice $T$ so that
$p^aT_0\subset T\subset T_0$. Take the schematic closure of
$T/p^aT_0$ in the finite flat group of $p^a$-torsion, and quotient the
good BT group by this finite flat subgroup. Its Tate lattice is
$p^{-a}T$, which is isomorphic to $T$. Finite flat quotients of BT
groups give BT groups, level by level. This argument uses only the
rational existence theorem at $p=2$, and does not assume an integral
$p=2$ equivalence in Kisin's Theorem~2.2.7.
\end{proof}

\begin{theorem}[B-RD17: a finite universal index exists]
\label{thm:B-RD17}
\status{Proved; independent proof and primary-source audit passed.}
For every prime $p$, every absolute ramification index $e$, and every
height $h\geq1$, a finite universal character-determination index in
B.4 exists. With $D_{p,h}$ as above, set
\[
 q=\left\lceil\log_p(eD_{p,h})\right\rceil,
 \qquad C=(2q+2)(q+2).
\]
An explicit sufficient index in the manuscript's convention is
\[
                       \nu=2Ch^2.                    \tag{B-RD17.4}
\]
In other words, the prime-to-$p$ manuscript character is determined
by $M[p^{\nu+1}]$ at this index. This bound is not asserted to be
sharp. For every prime and $h=1$, Proposition~\ref{prop:B-RD5} gives
the sharp value $\nu=0$; Proposition~\ref{prop:B-RD18} supplies
nonzero lower bounds for every admissible ramification in higher
height.
\end{theorem}
\begin{proof}
Choose $L$ as in Lemma~\ref{lem:B-RD17-degree}. Its absolute
ramification $e_L$ satisfies $e_L\leq eD_{p,h}$.
By \cite[Corollary~2 and the bound following Theorem~1, p.2 of the
author version]{VasiuZink2012Boundedness}, the generic Hom extension
exponent can be chosen to satisfy
\[
 c\leq(\log_p e_L+v_p(e_L)+2)(v_p(e_L)+2)
       \leq(2q+2)(q+2)=C.
\]
At level $n=2Ch^2+1$, Lemma~\ref{lem:B-RD17-common} applies.
Since $M(\nu)=M[p^{\nu+1}]$, this is precisely (B-RD17.4).
\end{proof}

\begin{remark}[The remaining sharp problem]
Theorem~\ref{thm:B-RD17} excludes the alternative that no finite
index exists. It does not determine the least index in height
$h\geq2$, and does not turn the manuscript monodromy condition into
potential semistability. The logarithmic cyclotomic extension in
Proposition~\ref{prop:B-RD4} remains a counterexample to that claim.
The refinement proof above works directly with good graded pieces
and therefore includes that extension and analogous nongood
extensions.
\end{remark}
